% Source: https://artofproblemsolving.com/community/c3h3176269_2023_mbmt_geometry_round__montgomery_blair_math_tournament
% Stable identity: aops-topic-3176269-q09
The triangle altitude is $H=\dfrac{s\sqrt3}{2}$, and the inradius is $\dfrac{H}{3}$. Each cut lies at height $\dfrac{2H}{3}$ above its parallel side, leaving a corner triangle of altitude $\dfrac{H}{3}$ and side $\dfrac{s}{3}$. Subtracting the three corner areas from the original gives central area $\dfrac{\sqrt3s^2}{4}-3\dfrac{\sqrt3s^2}{36}=\dfrac{\sqrt3s^2}{6}$. Equating this to $\dfrac{9\sqrt3}{2}$ gives $s^2=27$, hence $\boxed{s=3\sqrt3}$.
\begin{asy}
usepackage("amsmath");
size(7cm);
pair A=(0,0),Bp=(6,0),C=(3,3*sqrt(3)),O=(3,sqrt(3));draw(A--Bp--C--cycle);draw(circle(O,sqrt(3)));pair[] H={(2,0),(4,0),(5,sqrt(3)),(4,2*sqrt(3)),(2,2*sqrt(3)),(1,sqrt(3))};for(int j=0;j<6;++j)draw(H[j]--H[(j+1)%6],blue);label("$s$",(3,0),S);dot("incenter",O,S);
\end{asy}
